We study whether a recursive large language model (LLM) ecosystem has a minimum viable population of sufficiently distinct agents below which it collapses. Prior work on model collapse focuses on single recursive training loops, while the population-level question remains open and even the meaning of an ``individual'' agent is underspecified. We address this gap with a bounded recursive inheritance experiment that uses a frontier API model as the generator, behavioral distance to define individuality, and four generations of synthetic handoff under both replacement and accumulation regimes.

Our setup evaluates six conditions that vary population size, prompt-level role diversity, and access to a persistent human anchor. Across 576 recursive outputs, we find no clean threshold in which larger populations reliably avoid collapse. Instead, the dominant pattern is an immediate generation-1 truthfulness shock across every condition, with truth rates dropping from 0.67--0.92 at generation 0 to 0.42--0.50 at generation 1, followed by partial recovery to 0.67--0.83 by generation 3. Role-diverse agents are behaviorally distinct by our assay, with an individuality ratio of 1.1873 versus 0.9294 for homogeneous agents, but they do not protect the ecosystem. Final-generation collapse severity is positively correlated with the individuality proxy ($\rho=0.8281$, $p=0.0418$), and the best final condition is the single-agent accumulation regime.

These results argue against a universal minimum viable population threshold in this proxy setting. They instead point to inheritance quality and preserved human grounding as the more important determinants of recursive stability.
